\section{Problem Statement and Hypotheses}
\label{sec:problem}

\subsection{The MQAR task}
Following Zoology~\citep{arora2024zoology}, fix a vocabulary $V$ with a set of keys
$K \subset V$ and a disjoint set of values $U \subset V$. A context is a token sequence of
length $n$ encoding $D$ distinct bindings $(k_i, v_i)$, $i \in [D]$, followed by (or interleaved
with) $Q$ query positions; at a query position the input is a key bound earlier in the context
and the model must emit the associated value. Because the binding $k_i \mapsto v_i$ is resampled
per sequence, the answer cannot be memorized in the weights---it must be retrieved from context
at inference time. We write $\varepsilon$ for the expected fraction of query positions answered
incorrectly under the task distribution $\mathcal{D}$. For a recurrent model the resource of
interest is the \emph{state}: the fixed-width summary carried across positions, with size $S$
measured in bits. The central object of study is
\begin{equation}
  S^*(D, Q, \varepsilon; \mathcal{D}) \;=\; \text{minimum recurrent state (bits) to achieve query error} \le \varepsilon \text{ on } \mathcal{D}.
\end{equation}

\subsection{What is already proven (not our contribution)}
We restate the prior art precisely so as never to reclaim it. Any model depending causally on a
binary input and solving MQAR requires $\Omega(n)$ bits of recurrent state in the worst case
(BASED Thm.~3.1~\citep{arora2024based}) \textbf{[PROVEN]}; copying requires linearly growing
state~\citep{jelassi2024copying} \textbf{[PROVEN]}; and $o(n)$-bit RNNs fail in-context
retrieval even with chain-of-thought~\citep{wen2025rnns} \textbf{[PROVEN]}. A from-scratch
``recurrent models need $\Omega(n)$ state for recall'' theorem is therefore closed; our
contribution must lie in a different regime or in the mechanism.

\subsection{The open question}
Every bound above is worst-case over (near-)uniform inputs, where the context carries
$\Theta(n)$ bits. Structured streams have low realized entropy $H$: keys repeat, distributions
are skewed, and the number of novel bindings carried can be $H \ll n$.

\begin{quote}
\emph{Is the recall--state requirement governed by the entropy $H$ of the
key$\rightarrow$value map realized under $\mathcal{D}$, rather than by raw length $n$? And can a
recurrent mechanism allocate state proportional to $H$---matching attention's recall on
structured inputs at sub-$O(n^2)$ cost---while falling back to the $\Omega(n)$ worst case only
when $H = \Theta(n)$?}
\end{quote}

\subsection{Hypotheses}

\paragraph{H1 (entropy-parameterized lower bound) --- established as a proposition.}
The lower-bound half of H1 is proven in Section~\ref{sec:theory}: under the recurrent-bottleneck
assumption, Proposition~\ref{prop:floor} gives
$S \ge H(M) - D\,[H_b(\varepsilon) + \varepsilon \log_2(|U|-1)]$, which recovers the $\Omega(n)$
bound when the realized map entropy $H(M) = \Theta(n)$ and is far smaller when $H(M) \ll n$
\textbf{[PROVEN]}. As anticipated by our honesty flag, the proof is standard (data processing and
Fano), so we present it as a proposition rather than a deep theorem and \emph{retire} the earlier
guessed form $\Omega((1 - H_b(\varepsilon))H)$; the contribution is the entropy
\emph{reparameterization} of the known length-based bound, which isolates the low-entropy regime
where a recurrent mechanism is not information-theoretically barred from attention-level recall.
What a lower bound cannot decide is \emph{achievability}---whether a mechanism actually reaches
$S \approx \Theta(H(M))$---which is the substance of H2 and is settled only empirically.

\begin{conjecture}[H2: entropy-gated sparse memory]
A subquadratic backbone augmented with a small external associative memory whose \emph{writes}
are gated by a surprise/entropy signal (write only novel, high-surprise bindings) and whose
\emph{reads} use approximate nearest-key retrieval attains attention-level recall on structured
MQAR ($H \ll n$) at state $\approx \Theta(H)$, pushing the empirical recall--memory Pareto
frontier past matched-state-budget baselines, while matching them (no better, no worse) on
uniform inputs where $H = \Theta(n)$.
\end{conjecture}

\subsection{Falsification plan}
Proposition~\ref{prop:floor} holds under assumption (A2) (a fixed-width state through which all
answers factor); a model that re-reads the context, such as attention, is simply outside its
scope rather than a counterexample. The empirically falsifiable claim is therefore H2: it is
falsified if, on structured MQAR at matched (state, compute) budgets, the gated memory fails to
beat baselines on the accuracy-vs-resource frontier, or if any apparent gain vanishes under
properly tuned optimization (the caveat of \citet{okpekpe2025revisiting}). A rigorous negative
result is reported as a contribution. The empirical study (H2) is the primary deliverable and is
fully reproducible on the target hardware (a 4\,GB laptop GPU).
